\documentclass[UTF8,fontset=fandol,11pt,a4paper]{ctexart}
\usepackage[a4paper,margin=1.75cm,headheight=15pt]{geometry}
\usepackage{amsmath,amssymb,mathtools,booktabs,tabularx,array,enumitem}
\usepackage[most]{tcolorbox}
\usepackage{xcolor,fancyhdr,hyperref,microtype,lastpage}
\newcolumntype{Y}{>{\raggedright\arraybackslash}X}
\newcolumntype{L}[1]{>{\raggedright\arraybackslash}p{#1}}
\newcolumntype{C}[1]{>{\centering\arraybackslash}p{#1}}
\setlength{\parindent}{2em}
\setlength{\parskip}{0.34em}
\setlength{\emergencystretch}{3em}
\renewcommand{\arraystretch}{1.2}
\setlength{\tabcolsep}{3pt}
\setlist[itemize]{itemsep=0.18em,topsep=0.35em,leftmargin=2em}
\setlist[enumerate]{itemsep=0.18em,topsep=0.35em,leftmargin=2em}
\definecolor{deep}{RGB}{38,58,72}
\definecolor{soft}{RGB}{244,247,248}
\definecolor{greenish}{RGB}{52,91,74}
\definecolor{warnc}{RGB}{136,61,58}
\definecolor{purple}{RGB}{84,72,112}
\hypersetup{colorlinks=true,linkcolor=deep,urlcolor=purple,pdftitle={CO-01 数据表示与运算}}
\pagestyle{fancy}
\fancyhf{}
\lhead{\small\color{deep}\textbf{408 · 计算机组成原理 CO-01}}
\rhead{\small\color{deep}Bits · Meaning · Operation · Loss}
\cfoot{\small 第 \thepage\ 页 / 共 \pageref{LastPage} 页}
\renewcommand{\headrulewidth}{0.4pt}
\newtcolorbox{corebox}[1]{breakable,colback=soft,colframe=deep,title=\textbf{#1},boxrule=0.8pt,arc=2mm}
\newtcolorbox{methodbox}[1]{breakable,colback=white,colframe=greenish,title=\textbf{#1},boxrule=0.7pt,arc=1.5mm}
\newtcolorbox{warnbox}[1]{breakable,colback=white,colframe=warnc,title=\textbf{#1},boxrule=0.7pt,arc=1.5mm}
\newcommand{\kw}[1]{\textbf{\color{deep}#1}}
\title{\vspace{-1.1cm}\textbf{数据表示与运算}\\[0.4em]\large 有限位宽怎样保持可解释，并显式暴露信息丢失}
\author{408 · 计算机组成原理 Topic CO-01}
\date{}

\begin{document}
\maketitle

\begin{corebox}{母问题：同一串比特怎样得到确定的数学与硬件结果？}
比特没有自带 signed、float 或 address 语义。完整推理链是
\[
\begin{aligned}
&\text{Bit Pattern}\to\text{Interpretation}\to\text{Exact Value}\\
&\to\text{Operation}\to\text{Finite-width Result}\to\text{Lost Information / Flags}
\end{aligned}
\]
任何答案都应先固定 $(\text{width},\text{interpretation},\text{operation},\text{destination})$；溢出、舍入和异常都发生在“精确数学对象重新进入有限可表示集合”的边界。
\end{corebox}

\begin{methodbox}{本册学习范围}
\begin{itemize}
\item 本册解释位串如何表示整数、浮点数和校验信息，以及定长硬件如何完成加减乘除、移位和舍入。
\item 题目若给出结果位宽、舍入模式或状态位定义，以题设为准；本册先把这些条件翻译成可计算的规则。
\item 指令编码、完整控制通路和高级语言转换只在需要说明接口时出现；本册始终先把“位模式、数学值、目标格式”分开。
\end{itemize}
\end{methodbox}

\section{术语先行：读者先知道每个词在这里指什么}
\begin{methodbox}{五个基础词}
\begin{itemize}
\item \kw{位串（bit pattern）：}由 0 和 1 组成的固定长度序列；它本身没有正负号或小数意义。
\item \kw{位宽（width）：}位串包含的位数。硬件只保留这个宽度内的结果，超出的位必须被截掉或另行保存。
\item \kw{无符号/有符号：}同一位串的两种整数解释；无符号把每一位看成正的二进制权重，有符号补码把最高位看成负权。
\item \kw{ALU：}算术逻辑单元，负责加、减、与、或、异或、比较和移位等基本运算，并可产生进位、零、负和溢出等标志。
\item \kw{IEEE 754：}规定二进制浮点数的字段、特殊值和舍入规则的标准；本册把它当作一种格式，而不是某个处理器品牌。
\end{itemize}
\end{methodbox}

\tableofcontents
\newpage

\section{表示系统：位串、解释与可表示集合}
一个 $n$ 位位串只是集合 $\{0,1\}^n$ 的元素。unsigned 把它解释为
\[
U(b)=\sum_{i=0}^{n-1}b_i2^i,\qquad 0\le U\le 2^n-1;
\]
补码把最高位赋负权：
\[
T(b)=-b_{n-1}2^{n-1}+\sum_{i=0}^{n-2}b_i2^i,
\quad -2^{n-1}\le T\le2^{n-1}-1.
\]
同一位串可同时作为无符号数、有符号数、掩码、地址或浮点字段；意义来自消费它的操作，而不是某一位的外观。

原码把符号与绝对值分开，存在 $+0/-0$，加减需额外处理符号和大小。补码把定长加减统一到模 $2^n$ 环，只有一个零，代价是范围不对称：最小负数没有同宽正数。

\section{定长加减：同一低位结果，不同溢出解释}
$n$ 位加法器保留
\[
r=(a+b)\bmod 2^n.
\]
unsigned overflow 关注最高位 carry；signed overflow 关注精确结果是否落在补码范围。加法中“同号输入得到异号输出”可判 signed overflow；等价电路判据是符号位的进位与出位异或。

\begin{center}\small
\begin{tabularx}{\linewidth}{L{3.0cm}YY}\toprule
4 位运算 & unsigned 解释 & 补码解释\\\midrule
$1111+0001\to0000$ & $15+1$ 越界，carry=1 & $-1+1=0$，无 overflow\\
$0111+0001\to1000$ & $7+1=8$，无 unsigned 越界 & $7+1$ 越界，overflow=1\\
$1000$ 取负仍为 $1000$ & 位模式变换结果 & $-8$ 的同宽正数不可表示\\
\bottomrule
\end{tabularx}
\end{center}

减法先明确硬件采用 $a+(\overline b+1)$，再按题目对 carry/borrow 的定义解释。零、负、carry、overflow 是不同谓词；不能由一个标志推出另一个解释。

\section{扩展、截断与移位都在选择保留什么}
zero extension 保留 unsigned 值；sign extension 保留补码值。截断到低 $k$ 位相当于模 $2^k$，通常不保留数学值。判断窄化是否无损的通用办法是：截断后按原解释扩回，是否恢复原位串/数值。

逻辑移位补零；算术右移复制补码符号位，常对应向负无穷取整的除 $2^k$，不自动等于高级语言的有符号除法。左移低位结果与乘 $2^k$ 同余，但是否溢出仍取决于目标解释和被丢高位。

\begin{methodbox}{不变量}
每次扩展、截断或移位都问：要保留的是\kw{位模式、无符号值、补码值，还是仅保留模结果}？若答案不明确，后续标志判断没有语义基础。
\end{methodbox}

\section{乘法：完整乘积、部分积与目标宽度}
两个 $n$ 位无符号数的完整乘积最多需要 $2n$ 位；两个 $n$ 位补码数的完整乘积也用 $2n$ 位承载。若只保留低 $n$ 位，unsigned 与补码乘法的低位位串相同，但高位和溢出解释不同。

窄结果判据：unsigned 要求完整乘积高 $n$ 位全零；补码要求完整乘积高 $n$ 位等于低半部符号位的扩展。更一般地，先得到足宽精确积，再检查是否能无损压回目标格式。

\subsection{移位加法与原码一位乘}
无符号乘法把乘数每一位 $q_i$ 选择的 $M\cdot2^i$ 累加。序贯乘法器用部分积、乘数和被乘数移位，以时间换加法器/寄存器资源；并行阵列同时生成和压缩部分积，以面积与关键路径换吞吐。

原码一位乘先取符号异或，再对绝对值做无符号部分积；它不是补码 Booth 算法的同义词。题目给出联合寄存器时，必须先写每段宽度、移位方向和每拍不变量，不能跨教材套门数或节拍数。

\subsection{Booth 的生成逻辑}
Radix-2 Booth 用相邻位 $(Q_0,Q_{-1})$ 压缩连续的 1：$01$ 加 $M$，$10$ 减 $M$，$00/11$ 不变，然后对联合寄存器算术右移。它利用
\[
00111100_2=01000000_2-00000100_2
\]
把一串加法变为边界处加/减。若题目使用更高基 Booth、不同位对次序或不同寄存器布局，先声明版本；“Booth 对所有输入都更快”不是稳定结论。

\section{除法：商余关系是最终不变量}
除法的稳定正确性条件是
\[
X=YQ+R,\qquad |R|<|Y|,
\]
余数符号和商的取整方向服从题设/ISA。恢复余数法在试减为负时恢复；不恢复余数法让当前余数符号决定下一拍加还是减，省去立即恢复，但结束时可能需要校正。

算法题先固定 partial remainder、quotient/dividend、divisor 和附加位布局。每拍写判断位、加减对象、移位方向和上商位；最终用商余关系回代。原码与补码算法、左移余数与右移除数版本不能混用同一口诀。

除数为零可由零检测逻辑发现，但架构结果可能是同步异常、规定值或其他语义。补码最小负数除以 $-1$ 会产生超出同宽 signed 范围的商；“同位宽整数除法不溢出”是错误命题。

\section{浮点：在非均匀格点上表示实数近似}
普通二进制规格化数为
\[
x=(-1)^s(1.f)_2\,2^{E-Bias}.
\]
符号选方向，阶码选尺度，尾数选该尺度上的格点。阶码全 0 或全 1 必须分支处理 zero、subnormal、infinity 和 NaN，不能套普通隐藏位公式。

subnormal 使用 $0.f$ 和最小正常指数，提供渐进下溢；infinity 与 NaN 的传播、比较和异常细节以标准/题设为准。把整个 IEEE 754 编码称作“原码”会掩盖偏置阶码、隐藏位和特殊编码。

相邻可表示数的绝对间距随指数增长，因此“大数精度更高”与“有效位数固定”不是一回事：有效相对精度近似稳定，但绝对间距变大。多数十进制小数在二进制中无限展开，输入编码时已经发生第一次舍入。

\section{浮点运算链与舍入证据}
加减法的生成链是
\[
\text{Classify}\to\text{Align}\to\text{Significand Operation}\to\text{Normalize}\to\text{Round}\to\text{Encode/Signal}.
\]
对阶右移较小尾数，sticky 位应记录所有被移出低位的 OR；guard、round、sticky 与保留末位共同决定舍入。round-to-nearest, ties-to-even 下，恰好中点时让保留结果最低位为偶数；其他舍入模式必须另算。

乘除先处理符号和阶码，再算有效数，规格化、舍入并检查 exponent range。overflow 与 underflow 是目标格式边界；inexact 表示精确结果不在目标格点上。只看阶差达到尾数位数就断言小数完全无影响，可能漏掉 GRS 对最终舍入的作用。

\begin{warnbox}{三个浮点反例}
$0.1+0.2$ 涉及输入编码、运算和结果编码的多次有限映射；小量对阶后主体尾数为零，sticky 仍可能改变舍入；浮点加法不满足结合律，括号会改变中间舍入点。
\end{warnbox}

\section{五列概念边界}
\begin{center}\small
\begin{tabularx}{\linewidth}{L{1.8cm}C{0.35cm}L{1.9cm}YY}\toprule
概念 A & $\neq$ & 概念 B & 真正区别与题面信号 & 混淆后果\\\midrule
bit pattern & $\neq$ & value & 存储对象 vs 解释后的数学对象 & 同位串误判\\
carry & $\neq$ & overflow & unsigned 出位 vs signed 越界 & 标志互推\\
sign extension & $\neq$ & zero extension & 保补码值 vs 保 unsigned 值 & 扩宽变值\\
full product & $\neq$ & narrow result & 足宽精确积 vs 目标寄存器保留部分 & 漏乘法溢出\\
divide-by-zero detect & $\neq$ & ISA exception & 电路条件 vs 软件可见契约 & 从门电路推出 OS 行为\\
rounding & $\neq$ & truncation & 映射到最近/方向格点 vs 直接丢低位 & 浮点结果错一 ULP\\
\bottomrule
\end{tabularx}
\end{center}

\section{做题控制协议}
\begin{enumerate}
\item 写位宽、解释、运算和目标格式；
\item 先写可表示集合与精确数学结果；
\item 用足够中间位执行扩展、加减、部分积或商余更新；
\item 明确第一次丢信息的位置与舍入/截断规则；
\item 分开判断 carry、overflow、zero、sign、inexact 等状态；
\item 乘除用积回算或 $X=YQ+R$；浮点用分类、数量级和反向解码校验；
\item 最后才套 ISA 对标志、异常、饱和或返回值的契约。
\end{enumerate}

\begin{methodbox}{压缩成第一动作}
看到表示或运算题，先问：\kw{这串位按什么格式解释？精确结果在哪一步被压回多少位，丢掉了什么证据？}
\end{methodbox}

\section{本册与整机的接口}
\begin{corebox}{从位串到处理器状态}
本册把无语义的位串解释成给定位宽下的值、结果和信息损失证据。处理器后续会使用这些结果、宽度、符号扩展或舍入信息；某个标志是否成为软件可见异常，取决于具体指令集的规定。
\end{corebox}
本册的局部成本包括运算位宽、加法器/乘除迭代次数、关键路径、舍入逻辑和精度损失。它不直接给出程序总时间，性能题要把这些参数交给 `IC x CPI x clock period` 的全局成本口径。

\section{知识地图：从表示到可验证结果}
\begin{center}\small
\begin{tabularx}{\linewidth}{L{3.0cm}YY}\toprule
考点 & 本册机制 & 接口\\\midrule
整数表示与运算 & 位权、补码、扩展、标志与目标位宽 & CO-02/03\\
乘法器与 Booth & 部分积、联合寄存器、相邻位重编码 & CO-03 Use\\
除法与溢出 & 商余不变量、恢复/不恢复、边界检测 & ISA exception\\
IEEE 754 & 分类、对阶、规格化、GRS 舍入、特殊值 & ISA rounding/status\\
\bottomrule
\end{tabularx}
\end{center}

\begin{corebox}{一句话}
有限硬件的核心不是“把数学算错”，而是按明确格式把无限或足宽结果投影回有限集合；正确答案必须同时给出保留结果、丢失信息和其软件可见解释。
\end{corebox}

本册覆盖定点与浮点表示、ALU 加法器、乘除迭代、校验码和类型转换机制；大小端与对齐只保留表示层交接，完整地址语义由 CO-02 拥有。并行阵列只是“用更多硬件换更短时间”的一种实现；Booth 或不恢复除法的寄存器安排和拍数必须以题设为准。

\section{补充细节：从进制到可验证的位串状态机}

本节把所有表示与运算统一到同一条可检查链上。所有计算先经过
\[
\text{Representation}\to\text{Exact Operation}\to\text{Width Reduction}
\to\text{Status / Error Evidence}\to\text{ISA-visible Result}.
\]
其中每个箭头都对应一个可以被题目追问的状态变化：进制转换改变的是书写表示，补码运算改变的是定长环上的元素，ALU 产生的是低位结果和候选标志，浮点运算改变的是非均匀格点上的近似，校验码则改变合法码字集合的距离。

\subsection{进位计数制：先写位置计数式，再分离整数与小数}

任意 $r$ 进制数的展开式为
\[
x=\sum_{i=m}^{n}d_i r^i,\qquad 0\le d_i<r.
\]
这条式子同时处理整数位和小数位；题目中的“转换”只是选择不同的系数表示。$r$ 进制转十进制直接按权展开，十进制整数转 $r$ 进制反复除以 $r$ 取余并逆序，小数转 $r$ 进制反复乘以 $r$ 取整并正序。整数部分和小数部分不能混用同一方向的余数序列。

\begin{methodbox}{二进制与八/十六进制的分组理由}
因为 $8=2^3$、$16=2^4$，二进制小数点两侧分别从小数点向两端按 $3$ 位或 $4$ 位分组；左侧不足补高位零，右侧不足补低位零。补的是书写占位，不改变原数值。非 $2$ 的幂进制没有这种机械分组，必须回到位置计数式或乘除法。
\end{methodbox}

有限位二进制小数是否终止，取决于约分后的分母是否只含 $2$ 的因子。分母含 $5$ 的十进制小数通常会在二进制中无限展开，因此即使输入阶段没有显式算术，也已经发生一次舍入；后续浮点误差不能全部归咎于最后一步加法。

\subsection{真值、机器数与四种编码：同一位串必须先声明解释}

对 $n$ 位定点整数，真值是数学对象，机器数是位串。原码、反码、补码和移码分别使用不同的解释函数：
\[
\begin{aligned}
\text{原码: }&s\,|\,|x|,\\
\text{反码: }&x\ge0\Rightarrow x,\quad x<0\Rightarrow \text{正数位逐位取反},\\
\text{补码: }&x\ge0\Rightarrow x,\quad x<0\Rightarrow 2^n+x,\\
\text{移码: }&\text{把有符号值整体加固定偏置后编码}.
\end{aligned}
\]
原码与反码存在 $+0/-0$；补码只保留一个零，但多出一个最小负数；移码把大小关系转成无符号字典序，常用于阶码。任何“首位是 $1$ 所以为负数”的判断都必须先确认采用的是哪一种格式。

\begin{center}\small
\begin{tabularx}{\linewidth}{L{1.9cm}L{2.8cm}L{2.4cm}Y}
\toprule
编码 & 零的数量 & $n$ 位有符号范围（常见口径） & 运算/比较代价 \\
\midrule
原码 & $+0,-0$ & $-(2^{n-1}-1)\sim +(2^{n-1}-1)$ & 符号与数值分离，加减逻辑复杂 \\
反码 & $+0,-0$ & $-(2^{n-1}-1)\sim +(2^{n-1}-1)$ & 存在循环进位，仍需处理双零 \\
补码 & 一个零 & $-2^{n-1}\sim +(2^{n-1}-1)$ & 加减统一，溢出由目标宽度判断 \\
移码 & 取决于偏置约定 & 通常用于指数而非一般整数算术 & 无符号比较可对应有符号大小 \\
\bottomrule
\end{tabularx}
\end{center}

负数补码可用“按位取反加一”得到，但求相反数时必须连符号位一起处理；对最小负数求相反数仍得到自身的位串，原因是其正数不在同宽补码可表示集合内。转换题的安全流程是：先写真值，再写目标格式，最后检查是否超出目标范围；不要从一串位直接跳到十进制答案。

\subsection{补码加减、ALU 与标志：低位结果和解释必须分开}

在 $n$ 位补码环上，加法器实际计算
\[
S=(A+B)\bmod 2^n.
\]
减法通过
\[
A-B=A+\overline{B}+1
\]
复用同一加法器。低位 $S$ 是硬件保留的结果；它是否代表正确的 unsigned 或 signed 数，要另行判断。unsigned 的 carry-out 表示最高位之外产生了进位，减法中的 borrow 取决于具体加法器约定；signed overflow 的等价判据包括“同号相加得异号”或“最高位输入进位与输出进位不同”。

ALU 可以把加法、逻辑运算、移位和比较统一成“输入位串 + 控制码 -> 结果位串 + 状态标志”的函数。常见标志的职责不同：
\[
\begin{array}{lll}
C:&\text{无符号进位/借位证据},&
V:\text{有符号越界证据},\\
Z:&\text{结果全零},&
N/S:\text{结果最高位或符号证据}.
\end{array}
\]
题目若要求判断“是否溢出”，必须先写目标解释；不能把 $C=1$ 当作 $V=1$，也不能用结果符号位单独代替完整判据。

比较题进一步利用同一次减法的状态证据。令定长结果 $D=A-B$：$Z=1$ 可以判等；对补码有符号比较，若发生 overflow，结果最高位的表面符号需要被 $V$ 校正，因此常见 condition-code 模型用 $N\oplus V$（或 $SF\oplus OF$）表示“数学意义上的 $A<B$”。无符号大小关系则依赖 carry/borrow 的题设定义。若题目明确采用“减法时 $CF=1$ 表示借位”的口径，则 $A<B\Leftrightarrow CF=1$，$A>B\Leftrightarrow ZF=0\land CF=0$；有些 ISA 的 carry 位恰好采用“无借位”为 1，不能把这一口径跨题硬套。CO-01 只负责这些标志怎样从有限宽度结果产生，条件转移怎样消费它们由 CO-02/题设决定。

\begin{warnbox}{ALU 题的三层检查}
第一层在位串层检查每一位和最终进位；第二层按 unsigned 解释低位结果；第三层按 signed/补码解释并判断 $V$。例如两个正数相加得到最高位为 $1$，可能是 signed overflow，但这并不意味着 unsigned 结果非法。三层混成一步，最容易把 carry、borrow、overflow 写成同一个概念。
\end{warnbox}

\subsection{移位：方向、补入位与取整方向共同决定语义}

逻辑左移通常左侧丢弃、右侧补零；逻辑右移左侧补零；算术右移复制符号位，以保持补码符号。对无溢出的整数，左移 $k$ 位相当于乘 $2^k$，右移 $k$ 位只在题设允许的数值范围和取整口径下近似除以 $2^k$。补码算术右移对负数趋向负无穷，而“向零截断”可能需要额外修正，不能把两者混为同一除法。

循环移位把移出的位重新接到另一端，不发生数值意义上的补零；带进位循环移位还把进位标志作为额外一位状态。判断移位是否溢出，要比较被丢弃位是否是结果符号位的合法扩展，不能只看移位后低半部。

\subsection{定点乘法：先保足宽，再把结果压回目标格式}

无符号乘法的精确积最多需要 $2n$ 位；补码乘法也应先在足够宽度中保留完整乘积。若目标只保留 $n$ 位：
\begin{itemize}
  \item unsigned 无溢出要求被丢弃的高位全为 $0$；
  \item signed 无溢出要求被丢弃的高位全为保留结果符号位的扩展。
\end{itemize}
只观察低半部或某个最终 carry，不能证明乘法没有溢出。

Booth 算法把补码乘数的连续 $1$ 段改写为“加一次、减一次”的边界差。每一拍观察乘数当前最低位与附加位 $(Q_0,Q_{-1})$：$01$ 加被乘数，$10$ 减被乘数，$00/11$ 不加减；随后对联合的部分积、乘数和附加位做算术右移。附加位不是装饰，它保存了上一拍的边界；算术右移则保持部分积的补码符号。最后用足宽积或十进制乘法回算，检查位序和符号。

\subsection{迭代除法：恢复与不恢复共享商余不变量}

除法题先固定寄存器布局、被除数/除数符号、商位方向和循环次数。最终不变量是
\[
X=YQ+R,
\]
并且余数满足题设的符号和范围约束。恢复余数法在每一步左移部分余数并试减除数：若结果非负，上商 $1$；若结果为负，则恢复原余数并上商 $0$。它的优点是规则直观，代价是失败时多一次恢复加法。

不恢复余数法把“本拍余数为正/负”作为下一拍动作依据：正则减除数并上 $1$，负则加除数并上 $0$；最后若余数为负，再做一次修正。它省掉了每次立即恢复，但要求严格保持符号判断、左移顺序和末位修正。两种算法的逐拍表不能混用；题目给出寄存器图时，以图中的连接和移位方向为准。

\begin{methodbox}{乘除迭代题的最小逐拍表}
每拍至少记录五列：当前判断位/符号、加减对象、移位方向、上商或重编码位、循环后不变量。最后再做积回算或商余回算。若只能背“先移位还是先加减”的口诀而不能写出不变量，换一种寄存器布局就会失效。
\end{methodbox}

\subsection{IEEE 754：字段分配、特殊值和非均匀格点}

普通二进制规格化数写成
\[
x=(-1)^s(1.f)_2 2^{e-bias}.
\]
阶码使用偏置表示，便于比较和处理指数范围；规格化数隐藏最高有效位 $1$，把有限字段留给更多有效精度。字段全零或全一时必须进入特殊分支：
\[
\begin{array}{c|c|c}
E & F & \text{含义}\\\hline
0&0&+0/-0\\
0&\ne0&\text{subnormal}\\
\text{中间值}&\text{任意}&\text{normal}\\
\text{全一}&0&\pm\infty\\
\text{全一}&\ne0&\text{NaN}
\end{array}
\]

408 中最常用的两种二进制格式要把“字段位数”和“有效精度”分开记：
\begin{center}\small
\begin{tabularx}{\linewidth}{L{1.8cm}C{1.1cm}C{1.1cm}C{1.1cm}C{1.4cm}Y}
\toprule
格式 & $S$ & $E$ & $F$ & Bias & 规格化有效精度 \\
\midrule
binary32 / float & 1 & 8 & 23 & 127 & $1+23=24$ 位 \\
binary64 / double & 1 & 11 & 52 & 1023 & $1+52=53$ 位 \\
\bottomrule
\end{tabularx}
\end{center}
对 binary32，规格化数的阶码字段只使用 $E=1\sim254$，所以真实指数为 $-126\sim127$；$E=0$ 与 $E=255$ 留给零、非规格化数、无穷与 NaN 等特殊类。这里“阶码”是字段角色，“偏置表示/移码思想”是编码方法，两者不能写成同一个概念。binary64 同理，只是字段宽度和 Bias 改变。

subnormal 用固定的最小指数和 $0.f$ 提供渐进下溢；不能把它误套成隐藏位为 $1$ 的 normal。相邻浮点格点的绝对间距随指数增大，因而相对精度大致稳定而绝对间距变粗；“指数越大精度越高”与“有效位数固定”并不矛盾。

\subsection{浮点加减：对阶、规格化和 GRS 是连续状态}

两个有限浮点数相加的完整顺序为
\[
\text{Classify}\to\text{Align}\to\text{Significand Operation}
\to\text{Normalize}\to\text{Round}\to\text{Encode}.
\]
对阶时通常向较大阶码看齐，把较小尾数右移；右移丢掉的所有位先压成 sticky，再由 guard、round、sticky 和保留结果最低位决定舍入。相加可能产生最高位进位而右规，相减可能出现前导零而左规；两者都要同步修正阶码。舍入模式若是 round-to-nearest, ties-to-even，中点时保留结果最低位取偶数；题设改变舍入模式时不能继续套这个结论。

输入编码、有效数运算和结果编码都可能舍入，因此 $0.1+0.2$ 一类题不能只在十进制表面解释。浮点不满足结合律，是因为不同括号会改变中间的对阶、规格化和舍入位置。判断小量“完全无影响”也要看 sticky；即使主体尾数最终为零，被移出的低位仍可能改变最后一 ULP。

\subsection{校验码：把合法位串做成稀疏集合}

校验码的统一问题是：通信或存储发生错误后，接收方如何从冗余位判断“这串位是否仍是合法码字”，以及在能力范围内定位错误。码距 $d_{min}$ 是任意两个合法码字之间的最小汉明距离；检测 $t$ 位错误至少需要 $d_{min}\ge t+1$，纠正 $t$ 位错误至少需要 $d_{min}\ge2t+1$。

奇偶校验只给整个数据附加一个校验位，能检测奇数位翻转，不能定位错误，也不能可靠检测偶数位翻转。二维奇偶校验把行列约束叠加，可提供更强的定位能力，但具体能力仍由码的构造和错误模型决定。CRC 把位串看作二进制多项式，发送端在数据后补 $r$ 个零并除以生成多项式，余数作为校验序列；接收端对完整码字做同一除法，余数为零才通过。CRC 的“除法”是模 $2$ 加法，不是普通整数除法，减法等价于异或。

\subsection{海明码：综合校验结果就是错误位置}

若数据位数为 $k$、校验位数为 $r$，单错纠正要求
\[
2^r\ge k+r+1.
\]
校验位放在位置 $1,2,4,8,\ldots$，其余位置放数据。第 $j$ 个校验位负责所有位置编号的二进制表示中第 $j$ 位为 $1$ 的位置；编码时按偶校验填入，接收时重新计算各组校验结果并拼成 syndrome。syndrome 为 $0$ 表示没有检测到单错，非零值直接给出出错位编号。SEC-DED 再增加总体奇偶位：可实现单错纠正和双错检测，但三位及以上错误不应超出题设能力范围。

\subsection{C 类型转换：先保源值，再缩放位宽，最后按目标格式解释}

转换题最容易错在把 destination 的 signedness 提前用于 source。对 408 常见的二进制补码机器模型，可把整数内部转换压成三类：
\begin{itemize}
  \item \kw{长 $\to$ 短：}保留目标宽度的低位，高位被截断；新最高位随后按目标类型承担新的数值/符号角色；
  \item \kw{短 $\to$ 长：}先按\textbf{源类型}保值扩展。source unsigned 做零扩展，source signed 补码做符号扩展；
  \item \kw{同宽 signed $\leftrightarrow$ unsigned：}位模式不变，只改变解释函数。
\end{itemize}
因此 `signed short -> unsigned int` 不能看到 destination 是 unsigned 就立即补 0；应先把 signed short 的值符号扩展到目标位宽，再让同一目标宽度位串按 unsigned 重新解释。这个顺序可以压成
\[
\boxed{\text{Source Value}\to\text{Width Conversion}\to\text{Destination Interpretation}}.
\]

整型与浮点之间还要额外比较\textbf{范围}和\textbf{有效精度}。在常见 32 位 `int`、binary32 `float`、binary64 `double` 模型下：
\begin{center}\small
\begin{tabularx}{\linewidth}{L{2.4cm}YY}
\toprule
转换 & 稳定判断 & 主要风险 \\
\midrule
`int` $\to$ `float` & float 范围足够覆盖 32 位 int，但只有 24 位有效二进制精度 & 可能舍入，不是所有 int 都精确 \\
`int` $\to$ `double` & 53 位有效精度足以容纳 32 位 int & 可精确表示 \\
`float` $\to$ `double` & binary64 精度与范围都更大 & 有限 binary32 值可精确表示 \\
`double` $\to$ `float` & 精度和范围同时缩小 & 可能舍入，也可能溢出 \\
浮点 $\to$ 整数 & 目标格式没有小数部分 & 题设采用 C 常见口径且结果在目标范围内时向 0 截断；越界语义看语言/ISA/题设 \\
\bottomrule
\end{tabularx}
\end{center}
最稳的整数窄化检查仍是“截后再扩”：把结果按目标解释扩回，不能恢复原值就说明转换丢失了信息。高位重复并不自动等于安全，必须相对于 source、destination 和目标位宽判断。

\subsection{向 CO-02 交接：大小端与对齐不是第二套数值编码}
Endian 只规定多字节对象的字节怎样映射到递增地址，不改变寄存器中的数学值，也不反转单字节内部 bit；对齐规定对象合法/高效起始地址的边界。CO-01 只输出“对象位模式、宽度、扩展结果”，完整的地址—字节布局、结构体成员顺序、EA 与未对齐访问语义由 CO-02 拥有。

\subsection{补充细节到微架构骨架的回接表}

\begin{center}\small
\begin{tabularx}{\linewidth}{L{2.4cm}L{2.8cm}L{2.8cm}Y}
\toprule
题型入口 & 要固定的对象 & 硬件需要的结果 & 题目验证动作 \\
\midrule
进制/编码 & 位串、解释函数、可表示集合 & typed value、范围、比较顺序 & 先声明格式再读位权 \\
补码/ALU & 定长环、低位结果、$C/V/Z/N$ & 结果与候选状态位 & 分开 unsigned/signed 判定 \\
乘除迭代 & 部分积/余数、商、附加位 & 每拍写使能和移位需求 & 写循环不变量并回算 \\
浮点 & sign/exponent/fraction、特殊类 & 有效数、阶码、GRS、异常候选 & 分类后再对阶/舍入 \\
校验码 & 合法码字集合、syndrome & 错误检测/定位结果 & 用码距或多项式余数校验 \\
C 转换 & 源类型、目标类型、位模式 & 扩展/截断控制与访存宽度 & 截后再扩验证 \\
大小端/对齐 & 地址—字节映射与合法边界 & load/store 字节选择、事务拆分 & 画地址表，不凭字符串顺序猜 \\
\bottomrule
\end{tabularx}
\end{center}

\section{扩充后的自测与反向重建}

完成本册后，不能只背“公式出现在哪一节”，而应能从陌生题重新生成以下路径：
\begin{enumerate}
  \item 给出一串位，先判断题设格式、位宽和目标解释；
  \item 给出一个定点运算，保留足宽中间结果，分别判断 carry、overflow 和目标位宽截断；
  \item 给出一个乘除迭代寄存器图，写出每拍状态与最终回算不变量；
  \item 给出 IEEE 字段，先分类 zero/subnormal/normal/infinity/NaN，再进行运算或比较；
  \item 给出校验位布局，按码距、生成多项式或 syndrome 判断可检测/可纠正范围；
  \item 给出 C 转换、大小端或未对齐访问，明确位模式、地址字节和 ISA 行为是否改变；
  \item 最后把结果封装为 CO-02/CO-03 可使用的“值、宽度、标志、异常候选”状态包。
\end{enumerate}

\begin{methodbox}{本轮扩充的停止条件}
当读者能在不查表的情况下解释“为什么这个位串这样读、哪一步丢了信息、丢失是否属于溢出/舍入/校验失败、硬件下一步需要锁存什么”，CO-01 的机制层完成；具体题型速度、公式选择和考场退出仍属于计组 Rules 与题目验证，不继续把技巧堆入本册。
\end{methodbox}

\subsection{ALU 细节：从一位全加器到可控功能单元}
一位半加器的和与进位为
\[
S=A\oplus B,\qquad C=A\land B.
\]
一位全加器再接收输入进位 $C_{in}$：
\[
S=A\oplus B\oplus C_{in},\qquad
C_{out}=AB+AC_{in}+BC_{in}.
\]
把全加器串成 $n$ 位即可得到 ripple-carry adder；其延迟随进位逐位传播，最坏路径近似为 $O(n)$。先行进位把每位写成 propagate/generate：
\[
P_i=A_i\oplus B_i,\qquad G_i=A_iB_i,\qquad C_{i+1}=G_i+P_iC_i,
\]
并展开为
\[
C_1=G_0+P_0C_0,\quad C_2=G_1+P_1G_0+P_1P_0C_0,
\]
以更多并行逻辑和连线换更短的进位依赖。题目若比较串行进位与先行进位，比较的是微架构延迟、面积和扇入，不是改变 ISA 加法语义。

ALU 在加法器前可用 XOR/选择逻辑把 $B$ 取反并把 $C_0$ 置 1，从而复用同一硬件实现 $A-B$；在结果后接零检测、符号检测和溢出逻辑即可生成 $Z,N,V$ 等 flags。逻辑运算（AND/OR/XOR/NOT）、比较和移位可以由多路选择器接到同一结果总线；移位器是否在 ALU 内部由具体实现决定。控制字应说明功能选择、输入选择、进位输入和 flag 写使能，不能把“ALU 能算”误写成“所有 flags 每拍都更新”。

\subsection{CRC 与码距的可计算边界}
若数据多项式为 $M(x)$、生成多项式为 $G(x)$ 且 $G$ 的次数为 $r$，发送端先计算 $x^rM(x)$ 除以 $G(x)$ 的模 2 余数 $R(x)$，发送码字 $T(x)=x^rM(x)+R(x)$。接收端对 $T(x)$ 或收到的码字做同一模 2 除法，余数为零才通过。CRC 的“减法”是 XOR，不发生借位；生成多项式的首项和次数决定补零位数，不能用十进制除法替代。

若合法码字最小汉明距离为 $d_{min}$，检出至多 $e$ 位错误需 $d_{min}\ge e+1$，纠正至多 $t$ 位错误需 $d_{min}\ge2t+1$；若题目同时要求检出 $e$ 位并纠正 $t$ 位，常用充分条件为 $d_{min}\ge e+t+1$（$e\ge t$）。这些是码的结构能力，不等于某一次随机错误一定发生。

\begin{methodbox}{CO-01 的双重验算}
算术题用十进制/数学关系回算位串，校验题用码距/多项式余数回算编码，浮点题用分类—对阶—规格化—GRS—舍入回算。若两条独立路径不一致，优先检查目标位宽、符号扩展、补零方向和编址单位，而不是直接修改最后一位答案。
\end{methodbox}

\end{document}
